% ============================================================================
% Математика проективного многообразия (Книга 5)
% (book/sources/volumes/vol-projective-manifold-mathematics.tex)
% ============================================================================

\chapter*{Предисловие. Отделение чистой проективной математики от физики носителя}
\markboth{Предисловие}{Предисловие}
\addcontentsline{toc}{chapter}{Предисловие. Отделение чистой проективной математики от физики носителя}

Физика дискретного носителя (изложенная в \textbf{Книге~1}) и чистая математика на полярном фазовом кольце $\mathbb{Z}_{4Q}$ имеют различный предмет и различные критерии истинности. Физическая теория проверяется совпадением с константами природы и лабораторным экспериментом в чистой комнате. Математическая же теория представляет собой замкнутую дедуктивную систему тождеств и инвариантов над целыми числами $\mathbb{Z}$ и полями $\mathbb{Q}$, $\mathbb{R}$, $\mathbb{C}$, истинность которой устанавливается исключительно формальным выводом в ядре доказательства теорем.

Естественными координатами как в физике носителя, так и во всей настоящей математической монографии являются **целочисленные полярные координаты $(M, \varphi) \in \mathbb{N}^+ \times \mathbb{Z}_{4Q}$** --- радиус (амплитуда) $M \ge 1$ и полярный фазовый угол $\varphi \in \{0, 1, \dots, 4Q-1\}$. На ротаторе $|\mathcal{C}_Q(\varphi)| + |\mathcal{S}_Q(\varphi)| \equiv Q$ переход от полярных координат $(M, \varphi)$ к декартовым проекциям $(X, Y) = \bigl(M\mathcal{C}_Q(\varphi),\; M\mathcal{S}_Q(\varphi)\bigr)$ и обратный переход к полярному радиусу и углу
\[
  M \;=\; \frac{|X| + |Y|}{Q}, \qquad \varphi \;=\; \operatorname{atan2}_{L^1}(Y, X) \;=\; \frac{Q\,Y}{|X| + |Y|} \quad (0 \le \varphi \le Q)
\]
осуществляются **взаимно однозначно (биективно) в чистых целых числах** --- без извлечения квадратного корня $\sqrt{X^2+Y^2}$ и без бесконечных рядов Тейлора. Любые дробные отношения (тангенс $\mathcal{S}_Q/\mathcal{C}_Q$, точки единичной окружности Вейерштрасса $(X/R,\; Y/R)$, башни производных частных $D^n(f/g) = N_n/G_n$ и аппроксиманты Паде $P/Q$) хранятся парами целых чисел $\text{(числитель, знаменатель)}$ без преждевременного деления в плавающей точке и без потери остатка, представляя обычные числа из стандартных полей $\mathbb{Q} \subset \mathbb{R} \subset \mathbb{C}$.

В настоящую математическую монографию (\textbf{Книга~5}) собраны только чистые, конструктивно доказанные математические результаты (без неконструктивных аксиом выбора и без скрытых посылок в заголовках теорем):
\begin{enumerate}[label=\textbf{Глава \arabic*.},leftmargin=*]
\item \textbf{Дискретный ротатор $\mathbb{Z}_{4Q}$, проективные поля $\mathbb{Q}(i, \sqrt{d-1})$ и дифференциально-интегральное исчисление без делений:}
  кусочно-линейный характер Эйлера $(\mathcal{C}_Q, \mathcal{S}_Q)$ на алмазе $|\mathcal{C}_Q|+|\mathcal{S}_Q|=Q$, тождественно постоянный разностный вронскиан $\mathcal{W}_Q(\varphi)\equiv Q$, подъём Брахмагупты--Вейерштрасса на рациональную окружность $S^1(\mathbb{Q})$, отсутствие делителей нуля $2(\mathcal{C}_Q^2+\mathcal{S}_Q^2)\ge Q^2$, однородная арифметика без деления в полях $\mathbb{Q}(i)$, $\mathbb{Q}(\sqrt{d-1})$ и круговом поле $\mathbb{Q}(\zeta_8) = \mathbb{Q}(i,\sqrt{2})$, замкнутая башня производных частного $D^n(f/g) = N_n/G_n$ всех порядков $n\in\mathbb{N}$ с $G_n \ge 1$, точные дискретные интегралы Френеля, Якоби, Дирихле, Парсеваля и Бесселя, многомерная кубатура Симпсона на оболочке $O_2(d)$ и конструктивные когомологии де~Рама $d^2\equiv 0$.
\item \textbf{Изотропный дискриминантный конус в $\mathrm{SL}_2(\mathbb{Z})$, моноид башни аппроксимантов Паде и бескорневое кубическое сжатие:}
  матричное представление двухуровневых ступеней $M_{a,b}\in\mathrm{Mat}_{2\times 2}(\mathbb{Z})$ с инвариантным определителем $\det M_{a,b} = C$, изотропный конус $\frac{1}{4}((\operatorname{Tr}\widetilde{P})^2 - 4) = X^2 + Y^2 - T^2$, параболические жордановы матрицы $\widetilde{P}_k = \mp I_2 + k N_0 \in \mathrm{SL}_2(\mathbb{Z})$ ($N_0^2 = 0$), рекурсия Ганкеля--Вронскиана вдоль луча, мультипликативный моноид башни линейных аппроксимантов Паде $(P_{\mathcal{R}}, Q_{\mathcal{R}})$, бескорневое тождество кубического сжатия Паде третьего порядка с тождественным сокращением $X^3$ на произвольном этаже $\mathcal{R}$ и закон невырожденного ухода на соседний этаж башни Паде.
\item \textbf{Трёхшахтные планетарные дифференциалы $\mathcal{A} + \mathcal{B} = \mathcal{C}$, вронскиан Мейсона--Стотерса и проекторы Лагранжа--Вандермонда:}
  уравнение Виллиса $z_1(\omega_1-\omega_H)+z_2(\omega_2-\omega_H)=0$, ортогональный репер трёхшахтного дифференциала $\mathcal{F}_3$ ($\mathbf{r}\cdot\mathbf{\omega}=0$), факторизация определителя матрицы моментов Ганкеля $-\det\mathcal{H}_3 = \widetilde{A}\widetilde{B}\widetilde{C}(X-Y)^2(Y-Z)^2(Z-X)^2$, целочисленный кубический вронскиан Мейсона--Стотерса $W^* = A B^* - B A^*$ и определитель четвёртых степеней $-\det\mathcal{W}_{\mathrm{MS}} = d_{\mathrm{Bez}}(X^4+Y^4+Z^4)$, проекторы Лагранжа--Вандермонда по дефектам нижних этажей $(S_1, S_2, S_3)$, межэтажное тождество сцепления $k=2 \leftrightarrow k=3 \implies K=6$ рода $g(6)=10$ и стереографическое вложение любого трёхшахтного баланса в ротатор $\mathbb{Z}_{4Q}$.
\item \textbf{Алгебра многовинтовых планетарных каскадов и векторная суперпозиция встречных волн на сфере Римана $\mathbb{P}^1$:}
  двухдуговая композиция траекторий и хордовый вронскиан цикла, унимодулярный базис $\mathrm{SL}_2(\mathbb{Z})$ показателей $(u,v) = (2a-3b, 2b-a)$, встречное прямое-обратное сшивание Безу в $\mathbb{Z}_{P_1}\times\mathbb{Z}_{Q_2}$, моноид многовинтовых каскадов $(P_{L_1\mathbin{+\!+}L_2}, Q_{L_1\mathbin{+\!+}L_2}, D_{L_1\mathbin{+\!+}L_2})$, универсальная экстремальная вилка остатка по всем бинарным словам шестерён $3^B - 2^B \le R_w \le 2^{A-B}(3^B - 2^B)$ со строго положительной вариационной производной транспозиции соседних ступеней $\delta_{\mathrm{swap}}R_w = P_{L_1}Q_{L_2} \ge 1$, отражение граничных неподвижных точек $n_\infty^{(-)} = -1 \leftrightarrow n_\infty^{(+)} = 0 \leftrightarrow n^* = +1$, а также закон полного векторного сложения встречных волн в океане мод с тождественным сокращением чётной матрицы Грама $\mathrm{MainAmount}_{\mathrm{vec}}$, тройная дихотомия $M$-этажного каскада и конструктивная теорема о модельно-теоретической независимости.
\end{enumerate}

% ============================================================================
\chapter{Дискретный ротатор \texorpdfstring{$\mathbb{Z}_{4Q}$}{Z4Q}, однородные проективные поля \texorpdfstring{$\mathbb{Q}(i,\sqrt{d-1})$}{Q(i, sqrt(d-1))} и анализ без делений}
\label{ch:math-rotator-calculus}
% ============================================================================

\section{Характер Эйлера на \texorpdfstring{$L^1$}{L1}-алмазе, тождество постоянного вронскиана и подъём Брахмагупты--Вейерштрасса}

Пусть $Q \in \mathbb{N}^+$ ($Q \ge 1$) --- произвольный целочисленный четверть-период. На циклическом кольце $\mathbb{Z}_{4Q}$ из $4Q$ фазовых состояний чётная (косинусная) и нечётная (синусная) проекции $(\mathcal{C}_Q(\varphi), \mathcal{S}_Q(\varphi)) \in [-Q, +Q]^2$ определяются кусочно-линейным обходом четырёх квадрантов $L^1$-алмаза $|\mathcal{C}_Q(\varphi)| + |\mathcal{S}_Q(\varphi)| \equiv Q$: в первом квадранте $0 \le m \le Q$ имеем $(\mathcal{C}_Q(m), \mathcal{S}_Q(m)) = (Q - m, m)$ с продолжением на остальные квадранты поворотом на четверть периода $+Q$:
\begin{equation}
  \label{eq:math-rotator-quarter-shift}
  \mathcal{C}_Q(\varphi + Q) \;=\; -\mathcal{S}_Q(\varphi), \qquad
  \mathcal{S}_Q(\varphi + Q) \;=\; +\mathcal{C}_Q(\varphi).
\end{equation}

\begin{theorem}[Тождество постоянного разностного вронскиана, подъём Брахмагупты--Вейерштрасса, отсутствие делителей нуля и замкнутые дискретные квадратуры]
\label{thm:math-wronskian-weierstrass-quadratures}
\lean{thm1_constant_wronskian_weierstrass_lift_and_newton_leibniz, thm2_closed_form_fresnel_jacobi_dirichlet_parseval_bessel_integrals, thm34_rotator_projective_field_invertibility_zero_divisor_freedom_and_q_i_sqrt2_fields}
\leanok
Для любого четверть-периода $Q \ge 1$ на кольце $\mathbb{Z}_{4Q}$ выполняются следующие безусловные конструктивные тождества:
\begin{enumerate}[label=\textbf{(\arabic*)},leftmargin=*]
\item \textbf{Тождественное постоянство разностного вронскиана и формула Ньютона--Лейбница:}
  для первых разностей $\Delta\mathcal{C}_Q(\varphi) = \mathcal{C}_Q(\varphi+1) - \mathcal{C}_Q(\varphi)$ и $\Delta\mathcal{S}_Q(\varphi) = \mathcal{S}_Q(\varphi+1) - \mathcal{S}_Q(\varphi)$ дискретный вронскиан тождественно равен четверть-периоду $Q$ в каждой точке кольца (включая угловые вершины смены наклона), а конечный интеграл первой разности любой функции $F : \mathbb{Z} \to \mathbb{Z}$ телескопически точен:
  \begin{align}
    \mathcal{W}_Q(\varphi) &:= \mathcal{C}_Q(\varphi)\,\Delta\mathcal{S}_Q(\varphi) - \mathcal{S}_Q(\varphi)\,\Delta\mathcal{C}_Q(\varphi) \;\equiv\; Q, \label{eq:math-const-wronskian} \\
    \sum_{k=0}^{N-1} \bigl(F(a + k + 1) - F(a + k)\bigr) &\equiv F(a + N) - F(a). \label{eq:math-newton-leibniz}
  \end{align}
\item \textbf{Подъём Брахмагупты--Вейерштрасса на рациональную окружность $S^1(\mathbb{Q})$ и отсутствие делителей нуля на $L^1$-алмазе:}
  однородная тройка Вейерштрасса $(X_T, Y_T, R_T) = \bigl(\mathcal{C}_Q^2 - \mathcal{S}_Q^2,\ 2\mathcal{C}_Q\mathcal{S}_Q,\ \mathcal{C}_Q^2 + \mathcal{S}_Q^2\bigr)$ тождественно лежит на евклидовой окружности, а знаменатель $R_T = \mathcal{C}_Q^2 + \mathcal{S}_Q^2$ строго положителен во всех точках $L^1$-алмаза $|\mathcal{C}_Q| + |\mathcal{S}_Q| = Q \ge 1$:
  \begin{align}
    X_T^2 + Y_T^2 &= (\mathcal{C}_Q^2 - \mathcal{S}_Q^2)^2 + (2\mathcal{C}_Q\mathcal{S}_Q)^2 \;\equiv\; (\mathcal{C}_Q^2 + \mathcal{S}_Q^2)^2 \;=\; R_T^2, \label{eq:math-weierstrass-circle} \\
    2\,R_T &= 2\bigl(\mathcal{C}_Q^2 + \mathcal{S}_Q^2\bigr) \;\equiv\; (\mathcal{C}_Q + \mathcal{S}_Q)^2 + (\mathcal{C}_Q - \mathcal{S}_Q)^2 \;\ge\; Q^2 \;\ge\; 1. \label{eq:math-no-zero-divisors}
  \end{align}
\item \textbf{Однородная проективная обратимость в полях $\mathbb{Q}(i)$, $\mathbb{Q}(\sqrt{m})$ и круговом поле $\mathbb{Q}(i,\sqrt{2}) = \mathbb{Q}(\zeta_8)$:}
  в однородных проективных координатах без деления мультипликативные обратные в гауссовом поле $\mathbb{Q}(i)$ и квадратичном поле $\mathbb{Q}(\sqrt{m})$ ($m = d - 1$) задаются полиномиальными операторами
  \begin{align}
    \operatorname{gaussFieldInv}(X, Y, R) &:= \bigl(R X,\; -R Y,\; X^2 + Y^2\bigr), \label{eq:math-gauss-inv} \\
    \operatorname{quadFieldInv}_m(a, b, D) &:= \bigl(D a,\; -D b,\; a^2 - m b^2\bigr), \label{eq:math-quad-inv}
  \end{align}
  удовлетворяющими $(X, Y, R)\otimes_{\mathbb{Q}(i)}\operatorname{gaussFieldInv}(X, Y, R) = \bigl(R(X^2+Y^2),\, 0,\, R(X^2+Y^2)\bigr) \sim (1, 0, 1)$, а при $m = 2$ первообразный корень $8$-й степени из единицы $\zeta_8 = (0, 1, 0, 1, 2) \in \mathbb{Q}(i,\sqrt{2})$ удовлетворяет точным проективным соотношениям $\zeta_8^2 \sim i$, $\zeta_8^4 \sim -1$, $\zeta_8^8 \sim +1$, $\zeta_8\bar{\zeta}_8 \sim +1$.
\end{enumerate}
\end{theorem}

\begin{proof}
\leanok
В пункте~(1) внутри первого квадранта $(Q - m)(m + 1) - m(Q - m - 1) = Q m + Q - m^2 - m - Q m + m^2 + m = Q$, а в вершине $m = Q$ имеем $0\cdot(Q-1) - Q\cdot(-1) = Q$. В пункте~(2) раскрытие $(\mathcal{C}_Q^2 - \mathcal{S}_Q^2)^2 + 4\mathcal{C}_Q^2\mathcal{S}_Q^2 = \mathcal{C}_Q^4 + 2\mathcal{C}_Q^2\mathcal{S}_Q^2 + \mathcal{S}_Q^4 = (\mathcal{C}_Q^2 + \mathcal{S}_Q^2)^2$ и $(\mathcal{C}_Q + \mathcal{S}_Q)^2 + (\mathcal{C}_Q - \mathcal{S}_Q)^2 = 2(\mathcal{C}_Q^2 + \mathcal{S}_Q^2)$ при $|\mathcal{C}_Q| + |\mathcal{S}_Q| = Q$ даёт $2 R_T \ge Q^2 \ge 1$. В пункте~(3) перемножение однородных пятёрок в $\mathbb{Q}(i,\sqrt{2})$ проверяется прямым раскрытием скобок.
\end{proof}

\section{Замкнутая башня производных частного всех порядков \texorpdfstring{$D^n(f/g)$}{D\^{}n(f/g)}, многомерная кубатура Симпсона и когомологии де~Рама}

\begin{theorem}[Башня разностных производных частного $D^n(f/g) = N_n/G_n$ всех порядков $n\in\mathbb{N}$, вариационный принцип, многомерная кубатура Симпсона на $O_2(d)$ и конструктивные когомологии де~Рама]
\label{thm:math-c-infinity-simpson-derham}
\lean{thm3_continuum_integration_from_cuboctahedron_archimedes_layers, thm35_general_field_derivative_operator_quotient_wronskian_and_variational_euler_lagrange, thm38_c_infinity_higher_derivatives_taylor_newton_expansion_and_pole_free_quotients, thm5_constructive_de_rham_cohomology_nilpotent_algebra_and_banach_tarski_free_measure}
\uses{thm:math-wronskian-weierstrass-quadratures}
\leanok
Для любых целочисленных функций $f, g : \mathbb{Z} \to \mathbb{Z}$ с разностным оператором $(Df)(x) = f(x+1) - f(x)$:
\begin{enumerate}[label=\textbf{(\arabic*)},leftmargin=*]
\item \textbf{Рекуррентная башня производных отношения $D^n(f/g) = N_n(x)/G_n(x)$ во всех порядках $n\in\mathbb{N}$ без полюсов:}
  индуктивная пара числителя и знаменателя $n$-й производной частного
  \begin{align}
    N_0(x) &:= f(x), \qquad
    N_{n+1}(x) \;:=\; N_n(x+1)\,g(x) - N_n(x)\,g(x+n+1), \label{eq:math-quot-numer-tower} \\
    G_n(x) &:= \prod_{j=0}^n g(x+j) \label{eq:math-quot-denom-tower}
  \end{align}
  удовлетворяет точному шаговому тождеству $\bigl(N_n(x+1)G_n(x) - N_n(x)G_n(x+1)\bigr)g(x) \equiv N_{n+1}(x)G_n(x)$ и конечно-разностной формуле Ньютона по падающим факториалам, причём для знаменателя Вейерштрасса $g(t) = R_Q(t) = \mathcal{C}_Q(t)^2 + \mathcal{S}_Q(t)^2$ ($Q \ge 1$) при всех $n\in\mathbb{N}$ и всех $x\in\mathbb{Z}$ выполняется $G_n(x) \ge 1$.
\item \textbf{Многомерное разбиение на пояса и кубатура Симпсона на оболочке $O_2(d)$ при произвольной размерности $d \ge 2$:}
  сечение оболочки $O_2(d) = \{e \in \{-1,0,+1\}^d \mid \sum_{i=1}^d e_i^2 = 2\}$ из $2d(d-1)$ вершин вдоль любой координатной оси образует три пояса $\bigl(|Z_{-1}(d)|, |Z_0(d)|, |Z_{+1}(d)|\bigr) = \bigl(2(d-1),\ 2(d-1)(d-2),\ 2(d-1)\bigr)$, которые в сумме с центральным весом $w_0(d) = 2(d-1)(6-d)$ дают параболические веса Симпсона:
  \begin{equation}
    \label{eq:math-simpson-weights-d}
    \bigl(|Z_{-1}(d)|,\; |Z_0(d)| + w_0(d),\; |Z_{+1}(d)|\bigr) \;\equiv\; 2(d-1)\times(1,\; 4,\; 1),
  \end{equation}
  воспроизводя точный трёхмерный интеграл Римана $3\iiint_{[-1,1]^3} P\,\mathrm{d}V = 24c_0 + 8(c_x+c_y+c_z)$ на произвольном кубическом многочлене $P$, тогда как сумма по архимедовым слоям $S_k(3) = 10k^2 + 2$ удовлетворяет $3\bigl(1 + \sum_{k=1}^N (10k^2+2)\bigr) \equiv 10N^3 + 15N^2 + 11N + 3$.
\item \textbf{Конструктивные когомологии де~Рама $d^2 \equiv 0$ и обращение нильпотентных операторов в $\mathbb{Z}$:}
  кограничные операторы $d_0$ (на рёбрах), $d_1$ (на гранях) и $d_2$ (на кубах) удовлетворяют $d_1 \circ d_0 \equiv 0$ и $d_2 \circ d_1 \equiv 0$, циркуляция главных остатков по модулю $M$ вокруг замкнутой грани строго квантована периодом ($(d_1 r)_f = M\cdot n_f \in M\mathbb{Z}$), а любой нильпотентный оператор $\mathcal{M} = I + N_{\mathrm{nil}}$ ($N_{\mathrm{nil}}^3 = 0$) обращается в целых числах по формуле $(I + N_{\mathrm{nil}})^{-1} = I - N_{\mathrm{nil}} + N_{\mathrm{nil}}^2$.
\end{enumerate}
\end{theorem}

\begin{proof}
\leanok
Из определения~\eqref{eq:math-quot-denom-tower} следует $G_n(x+1)g(x) = g(x+n+1)G_n(x) = G_{n+1}(x)$, откуда вынесение общего множителя $G_n(x)$ даёт $(N_n(x+1)G_n(x) - N_n(x)G_n(x+1))g(x) = N_{n+1}(x)G_n(x)$. Тождество~\eqref{eq:math-simpson-weights-d} следует из $2(d-1)(d-2) + 2(d-1)(6-d) = 8(d-1)$, а равенства $d_1(d_0\varphi) \equiv 0$, $d_2(d_1 e) \equiv 0$ и $(I + N_{\mathrm{nil}})(I - N_{\mathrm{nil}} + N_{\mathrm{nil}}^2) = I + N_{\mathrm{nil}}^3 = I$ проверяются раскрытием скобок.
\end{proof}

\section{Биекция полярных и декартовых координат \texorpdfstring{$(M,\varphi)\leftrightarrow(X,Y)$}{(M, phi) <-> (X, Y)} без корней, точное обращение \texorpdfstring{$\operatorname{atan}_{L^1},\,\operatorname{acot}_{L^1}$}{atanL1, acotL1} и гиперболическое исчисление Гудермана--Минковского}

В естественных полярных координатах $(M, \varphi) \in \mathbb{N}^+ \times \mathbb{Z}_{4Q}$ на ротаторе $|\mathcal{C}_Q(\varphi)| + |\mathcal{S}_Q(\varphi)| \equiv Q$ переход к декартовым проекциям $(X, Y)$ и обратное восстановление полярного радиуса $M$ и угла $\varphi$, а также гиперболические функции и их обращение выполняются в чистых целых числах без бесконечных рядов, квадратных корней и логарифмов.

\begin{theorem}[Точная целочисленная биекция полярных и декартовых координат $(M,\varphi)\leftrightarrow(X,Y)$, обращение тангенса и котангенса, гиперболическая тройка Гудермана--Минковского и бескорневое обращение гиперболических функций]
\label{thm:math-inverse-trig-hyperbolic}
\lean{minkowski_hyperbola_poly_identity, exact_inverse_trig_and_rootfree_hyperbolic_calculus}
\uses{thm:math-wronskian-weierstrass-quadratures}
\leanok
Для любого четверть-периода $Q \ge 1$, произвольного полярного радиуса $M \ge 1$, произвольного полярного угла $\varphi \in \mathbb{Z}$ и координат светового конуса $U, V \in \mathbb{Z}$ выполняются следующие безусловные тождества:
\begin{enumerate}[label=\textbf{(\arabic*)},leftmargin=*]
\item \textbf{Точная целочисленная биекция полярных и декартовых координат $(M, \varphi) \leftrightarrow (X, Y)$ и обращение $\operatorname{atan}_{L^1},\,\operatorname{acot}_{L^1}$ без корней и рядов:}
  для декартовых проекций $X = M\,\mathcal{C}_Q(\varphi)$ и $Y = M\,\mathcal{S}_Q(\varphi)$ сумма модулей тождественно равна $|X| + |Y| = M\,Q$, поэтому полярный радиус $M$ и полярный угол $\varphi$ (вместе с дополнительным углом котангенса $Q - \varphi$ в первом квадранте $0 \le \varphi \le Q$) восстанавливаются с нулевой погрешностью при любом радиусе $M \ge 1$:
  \begin{equation}
    \label{eq:math-atan-acot-l1}
    M \;=\; \frac{|X| + |Y|}{Q}, \qquad
    \operatorname{atan}_{L^1}(Y, X) \;:=\; \frac{Q\,Y}{|X| + |Y|} \;\equiv\; \varphi, \qquad
    \operatorname{acot}_{L^1}(X, Y) \;:=\; \frac{Q\,X}{|X| + |Y|} \;\equiv\; Q - \varphi,
  \end{equation}
  причём $\operatorname{atan}_{L^1}(Y, X) + \operatorname{acot}_{L^1}(X, Y) \equiv Q$, а четырёхквадрантные операторы $\operatorname{cartToPolar}_Q(X, Y) = \bigl((|X|+|Y|)/Q,\ \operatorname{atan2}_{L^1}(Y, X)\bigr)$ и $\operatorname{acot2}_{L^1}(X, Y)$ биективно восстанавливают $(M, \varphi)$ во всех $4Q$ точках кольца $\mathbb{Z}_{4Q}$.
\item \textbf{Проективная гиперболическая тройка Гудермана--Минковского $(\mathcal{Ch}_Q, \mathcal{Sh}_Q, D_{\mathrm{hyp}})$ и инвариант светового конуса:}
  перестановка горизонтального катета и гипотенузы подъёма Вейерштрасса определяет целочисленные гиперболический косинус, гиперболический синус и гиперболический знаменатель:
  \begin{equation}
    \label{eq:math-gudermann-hyperbolic-triple}
    \mathcal{Ch}_Q(\varphi) \;:=\; \mathcal{C}_Q(\varphi)^2 + \mathcal{S}_Q(\varphi)^2, \qquad
    \mathcal{Sh}_Q(\varphi) \;:=\; 2\,\mathcal{C}_Q(\varphi)\,\mathcal{S}_Q(\varphi), \qquad
    D_{\mathrm{hyp}}(\varphi) \;:=\; \mathcal{C}_Q(\varphi)^2 - \mathcal{S}_Q(\varphi)^2,
  \end{equation}
  которые тождественно лежат на гиперболе Минковского и сохраняют сумму светового конуса по всему кольцу $\mathbb{Z}_{4Q}$:
  \begin{align}
    \mathcal{Ch}_Q(\varphi)^2 - \mathcal{Sh}_Q(\varphi)^2 &\equiv D_{\mathrm{hyp}}(\varphi)^2, \label{eq:math-minkowski-hyperbola-id} \\
    \mathcal{Ch}_Q(\varphi) + \bigl|\mathcal{Sh}_Q(\varphi)\bigr| &\equiv \bigl(|\mathcal{C}_Q(\varphi)| + |\mathcal{S}_Q(\varphi)|\bigr)^2 \;\equiv\; Q^2. \label{eq:math-hyperbolic-lightcone-const}
  \end{align}
\item \textbf{Обращение гиперболических функций без корней и без логарифмов ($\operatorname{invHyp}_Q$, $\operatorname{atanh}_{\mathrm{Cayley}}$, $\operatorname{acoth}_{\mathrm{Cayley}}$):}
  поскольку для любых $Q, \varphi \in \mathbb{Z}$ выполняется точная факторизация
  \begin{equation}
    \label{eq:math-invhyp-factorization}
    \mathcal{Ch}_Q(\varphi) + \mathcal{Sh}_Q(\varphi) - D_{\mathrm{hyp}}(\varphi) \;\equiv\; 2\,\mathcal{S}_Q(\varphi)\,\bigl(\mathcal{C}_Q(\varphi) + \mathcal{S}_Q(\varphi)\bigr),
  \end{equation}
  в первом квадранте $0 \le \varphi \le Q$ ($\mathcal{C}_Q(\varphi) + \mathcal{S}_Q(\varphi) = Q$) оператор обратного гиперболического синуса/тангенса восстанавливает точный фазовый угол $\varphi$ одним делением на $2Q$ без извлечения корней и без логарифмов:
  \begin{equation}
    \label{eq:math-invhyp-exact}
    \operatorname{invHyp}_Q\bigl(\mathcal{Ch}_Q(\varphi),\, \mathcal{Sh}_Q(\varphi),\, D_{\mathrm{hyp}}(\varphi)\bigr) \;:=\; \frac{\mathcal{Ch}_Q(\varphi) + \mathcal{Sh}_Q(\varphi) - D_{\mathrm{hyp}}(\varphi)}{2\,Q} \;\equiv\; \varphi,
  \end{equation}
  а на координатах светового конуса $(U, V)$ проективные гиперболический тангенс $\tanh_{\mathrm{Cayley}}(U, V) = (U - V : U + V)$ и котангенс $\coth_{\mathrm{Cayley}}(U, V) = (U + V : U - V)$ обращаются линейным оператором Кэли $\operatorname{atanh}_{\mathrm{Cayley}}(S_h, C_h) = (C_h + S_h,\ C_h - S_h) \equiv (2U, 2V)$.
\end{enumerate}
\end{theorem}

\begin{proof}
\leanok
В пункте~(1) при $X = \lambda(Q - \varphi) \ge 0$ и $Y = \lambda\varphi \ge 0$ имеем $|X| + |Y| = \lambda Q$, откуда $Q Y / (\lambda Q) = \varphi$ и $Q X / (\lambda Q) = Q - \varphi$. В пункте~(2) раскрытие скобок даёт $(c^2+s^2)^2 - (2cs)^2 = c^4 - 2c^2s^2 + s^4 = (c^2-s^2)^2$ и $c^2 + s^2 + 2|c||s| = (|c|+|s|)^2 = Q^2$. В пункте~(3) $(c^2+s^2) + 2cs - (c^2-s^2) = 2s^2 + 2cs = 2s(c+s) = 2Q\varphi$, а $(U+V \pm (U-V)) = (2U, 2V)$.
\end{proof}

\section{Конструктивный переход от квадратичных сумм на кольце \texorpdfstring{$\mathbb{Z}_{4Q}$}{Z\_4Q} к интегралам Френеля на \texorpdfstring{$\mathbb{R}$ ($\mathbb{C}$)}{R (C)} без рядов Тейлора и без таблиц}

Классическое значение интегралов Френеля $\int_0^\infty \cos(x^2)\,dx = \int_0^\infty \sin(x^2)\,dx = \frac{1}{2}\sqrt{\frac{\pi}{2}}$ выводится из дискретной квадратичной суммы на кольце $\mathbb{Z}_{512}$ чисто геометрически --- через $2$-адическое разложение кольца на пять оболочек, тождество Гаусса--Планшереля и спрямление многоугольника Вейерштрасса, без использования бесконечных степенных рядов Тейлора и таблиц:

\begin{theorem}[Общий символьный закон стационарной фазы Френеля для произвольного масштаба $b \in \mathbb{Z}$, тождество нормы Гаусса--Планшереля и безтейлоровский переход к $\int_0^\infty \cos(x^2)\,dx = \int_0^\infty \sin(x^2)\,dx = \frac{1}{2}\sqrt{\frac{\pi}{2}}$]
\label{thm:math-fresnel-cornu-continuum-bridge}
\lean{thm40_table_free_continuum_fresnel_cornu_brackets_and_gauss_sum_bridge}
\uses{thm:math-wronskian-weierstrass-quadratures, thm:math-c-infinity-simpson-derham}
\leanok
Для произвольного целочисленного масштаба $b \in \mathbb{Z}$ и фазового кольца размера $N_{\mathrm{ring}}(b) = 8b^2$ с четверть-периодом $Q(b) = 2b^2$, полупериодом $\Phi_{\mathrm{half}}(b) = 4b^2$ и точкой биссектрисы первого квадранта $K_P(b) = Q(b)/2 = b^2$ (при $b = 8$: $Q = 128,\ \Phi_{\mathrm{half}} = 256,\ N_{\mathrm{ring}} = 512,\ K_P = 64$) выполняются следующие безусловные символьные тождества:
\begin{enumerate}[label=\textbf{(\arabic*)},leftmargin=*]
\item \textbf{Общий закон попадания выживающей $2$-адической оболочки $\varphi(b, k) = b(2k+1)$ в биссектрису $45^\circ$ ($K_P(b) = b^2$) для всех $b, k, m \in \mathbb{Z}$:}
  на всех небиссекторных $2$-адических оболочках кольца квадратичные образы $\varphi^2 \bmod N_{\mathrm{ring}}(b)$ инвариантны относительно полупериодного сдвига $+\Phi_{\mathrm{half}}(b)$ и взаимно аннулируются по тождеству $\mathcal{E}(\psi + \Phi_{\mathrm{half}}(b)) = -\mathcal{E}(\psi)$, тогда как для любого элемента выживающей оболочки $\varphi(b, k) = b(2k+1)$ ($k \in \mathbb{Z}$) выполняется:
  \begin{equation}
    \label{eq:math-fresnel-four-shells-zero}
    \varphi(b, k)^2 - K_P(b) \;\equiv\; \Phi_{\mathrm{half}}(b)\,k(k+1),
  \end{equation}
  откуда при всех чётных $k = 2m$ и нечётных $k = 2m+1$ квадрат тождественно попадает в точку биссектрисы $K_P(b) = b^2 \leftrightarrow \pi/4$ по модулю полного периода кольца $N_{\mathrm{ring}}(b) = 8b^2$ (где $\mathcal{C}_{Q(b)}(K_P(b)) = Q(b) - K_P(b) = K_P(b) = b^2 = \mathcal{S}_{Q(b)}(K_P(b))$):
  \begin{align}
    \varphi(b, 2m)^2 &\equiv K_P(b) + N_{\mathrm{ring}}(b)\,m(2m+1), \label{eq:math-discrete-fresnel-even} \\
    \varphi(b, 2m+1)^2 &\equiv K_P(b) + N_{\mathrm{ring}}(b)\,(m+1)(2m+1), \label{eq:math-discrete-fresnel-odd} \\
    \mathcal{C}_{\mathrm{Fresnel}}(b) &= \mathcal{S}_{\mathrm{Fresnel}}(b) \;:=\; (4b)\,K_P(b) \;=\; 4\,b^3 \qquad \bigl(\text{при } b = 8:\; 32 \times 64 = 2\,048\bigr). \label{eq:math-discrete-fresnel-symmetry}
  \end{align}
\item \textbf{Общее символьное тождество нормы Гаусса--Планшереля для произвольного масштаба $b \in \mathbb{Z}$:}
  для любого $b \in \mathbb{Z}$ квадраты дискретных интегралов Френеля связаны с полупериодом $\Phi_{\mathrm{half}}(b) = 4b^2$, полным периодом кольца $N_{\mathrm{ring}}(b) = 8b^2$ и квадратом амплитуды $Q(b)^2 = 4b^4$ точными полиномиальными тождествами:
  \begin{align}
    \mathcal{C}_{\mathrm{Fresnel}}(b)^2 \;=\; \mathcal{S}_{\mathrm{Fresnel}}(b)^2 &\equiv \Phi_{\mathrm{half}}(b)\,Q(b)^2 \;=\; \frac{N_{\mathrm{ring}}(b)}{2}\,Q(b)^2 \;=\; 16\,b^6, \label{eq:math-fresnel-halfperiod-norm} \\
    \mathcal{C}_{\mathrm{Fresnel}}(b)^2 + \mathcal{S}_{\mathrm{Fresnel}}(b)^2 &\equiv N_{\mathrm{ring}}(b)\,Q(b)^2 \;=\; 32\,b^6. \label{eq:math-fresnel-plancherel-norm}
  \end{align}
\item \textbf{Прямой переход к $\int_0^\infty \cos(x^2)\,dx = \int_0^\infty \sin(x^2)\,dx = \sqrt{\pi/8}$ через спрямление многоугольника Вейерштрасса без рядов Тейлора:}
  при нормировке на амплитуду $Q(b)$ и $d_{\mathrm{in}}(\square) = 4$ квадранта кольца квадрат однолучевого интеграла Френеля равен $\bigl(\frac{\mathcal{C}_{\mathrm{Fresnel}}(b)}{4\,Q(b)}\bigr)^2 = \frac{N_{\mathrm{ring}}(b)}{2\cdot 4^2} = \frac{N_{\mathrm{ring}}(b)}{16}$. При спрямлении многоугольника Вейерштрасса с шагом $(\Delta t)^2 = 2\pi / N_{\mathrm{ring}}(b)$ множитель $N_{\mathrm{ring}}(b)$ сокращается тождественно при любом масштабе $b$:
  \begin{equation}
    \label{eq:math-fresnel-pi-over-8}
    \left(\frac{\mathcal{C}_{\mathrm{Fresnel}}(b)}{4\,Q(b)}\,\Delta t\right)^2
    \;=\; \frac{N_{\mathrm{ring}}(b)}{16}\cdot\frac{2\pi}{N_{\mathrm{ring}}(b)}
    \;\equiv\; \frac{\pi}{8}
    \quad\Longrightarrow\quad
    \int_0^\infty \cos(x^2)\,dx \;=\; \int_0^\infty \sin(x^2)\,dx \;=\; \sqrt{\frac{\pi}{8}} \;=\; \frac{1}{2}\sqrt{\frac{\pi}{2}},
  \end{equation}
  где $\pi \in [\pi_{\mathrm{in}}(n),\ \pi_{\mathrm{out}}(n)]$ зажимается рациональными суммами вписанного и описанного $n$-угольников Вейерштрасса~\eqref{eq:math-weierstrass-circle}, а квадратный корень извлекается рациональной башней Паде--$\mathrm{SL}_2(\mathbb{Z})$ Главы~\ref{ch:math-pade-sl2-towers} без рядов Тейлора и таблиц.
\end{enumerate}
\end{theorem}

\begin{proof}
\leanok
Раскрытие скобок даёт $(b(2k+1))^2 - b^2 = 4b^2 k(k+1) = \Phi_{\mathrm{half}}(b)k(k+1)$, а при $k = 2m$ и $k = 2m+1$ имеем $4b^2(2m)(2m+1) = 8b^2 m(2m+1)$ и $4b^2(2m+1)(2m+2) = 8b^2(m+1)(2m+1)$. Тождества~\eqref{eq:math-fresnel-halfperiod-norm}--\eqref{eq:math-fresnel-plancherel-norm} следуют из $(4b^3)^2 = 16b^6 = (4b^2)(2b^2)^2$.
\end{proof}

% ============================================================================
\chapter{Изотропный дискриминантный конус в \texorpdfstring{$\mathrm{SL}_2(\mathbb{Z})$}{SL2(Z)}, моноид башни аппроксимантов Паде и бескорневое кубическое сжатие}
\label{ch:math-pade-sl2-towers}
% ============================================================================

\section{Матричное представление ступеней, изотропный конус и параболические жордановы матрицы в \texorpdfstring{$\mathrm{SL}_2(\mathbb{Z})$}{SL2(Z)}}

Рассмотрим двумерную целочисленную решётку ступеней, порождённую парой величин $(U, V) = (A x^m, B y^n)$ с разностью $C = U - V$. Для любой внутренней ступени $(a, b)$ ($1 \le a \le m-1,\ 1 \le b \le n-1$) матрица ступени
\begin{equation}
  \label{eq:math-stage-matrix}
  M_{a,b} \;=\; \begin{pmatrix} A\,x^{m-a} & B\,y^{n-b} \\ y^b & x^a \end{pmatrix} \in \mathrm{Mat}_{2\times 2}(\mathbb{Z}),
  \qquad \det M_{a,b} \;=\; A\,x^m - B\,y^n \;=\; C
\end{equation}
определяет матрицу межступенчатого перехода $P = M_{a,b}^{\mathrm{adj}} M_{a',b'}$ с инвариантным определителем $\det P = C^2$ и нормированную унимодулярную матрицу $\widetilde{P} = \frac{1}{C}P \in \mathrm{SL}_2(\mathbb{Q})$.

\begin{theorem}[Изотропный дискриминантный конус $\mathrm{SL}_2(\mathbb{Z})$, параболические жордановы матрицы, рекурсия Ганкеля и счётная башня аппроксимантов Паде]
\label{thm:math-pade-tower-cubic-compression}
\lean{thm6_pillai_skeletal_matrix_decomposition_and_parabolic_sl2_jordan_triptych, thm7_pillai_recursive_descent_hankel_wronskian_farey_bezout_and_pade_thue_tower, thm10_root_free_cubic_pade_compression, thm10b_pillai_same_cell_two_solution_collision_and_gap_law, thm10c_pillai_arbitrary_floor_closed_cubic_residual_and_annihilation_condition, thm11_cross_floor_wronskian_and_p3_kernel}
\uses{thm:math-wronskian-weierstrass-quadratures}
\leanok
Для любых целых параметров $(U, V)$ с разностью $C = U - V$ и произвольного списка этажей башни Паде $\mathcal{R} = [r_1, \dots, r_M] \in \mathrm{List}\,\mathbb{Z}$ с индуктивным шагом $(P, Q) \mapsto \bigl(P + (r - 1)Q,\ (r - 1)P + Q\bigr)$ выполняются следующие безусловные тождества:
\begin{enumerate}[label=\textbf{(\arabic*)},leftmargin=*]
\item \textbf{Изотропный дискриминантный конус и нильпотентные жордановы образующие:}
  в координатах $X = (\widetilde{P}_{11} - \widetilde{P}_{22})/2$, $Y = (\widetilde{P}_{12} + \widetilde{P}_{21})/2$, $T = (\widetilde{P}_{12} - \widetilde{P}_{21})/2$ условие унимодулярности $\det\widetilde{P} = 1$ тождественно эквивалентно уравнению изотропного конуса
  \begin{equation}
    \label{eq:math-sl2-isotropic-cone}
    \frac{1}{4}\Bigl((\operatorname{Tr}\widetilde{P})^2 - 4\Bigr) \;=\; X^2 + Y^2 - T^2,
  \end{equation}
  а вдоль параболической образующей $\operatorname{Tr}\widetilde{P} = \mp 2$ ($X^2 + Y^2 = T^2$) семейство матриц $\widetilde{P}_k = \mp I_2 + k N_0$ с нильпотентной матрицей $N_0 = \begin{pmatrix} -1 & -1 \\ 1 & 1 \end{pmatrix}$ ($N_0^2 = 0$) принадлежит $\mathrm{SL}_2(\mathbb{Z})$ при всех $k \in \mathbb{Z}$.
\item \textbf{Гомоморфизм моноида башни Паде и точная факторизация суммы и разности на произвольном этаже $\mathcal{R}$:}
  для любых списков этажей $\mathcal{R}_1, \mathcal{R}_2 \in \mathrm{List}\,\mathbb{Z}$ композиция башни ассоциативна:
  \begin{align}
    P_{\mathcal{R}_1\mathbin{+\!+}\mathcal{R}_2}(U, V) &= P_{\mathcal{R}_2}\bigl(P_{\mathcal{R}_1}(U, V),\; Q_{\mathcal{R}_1}(U, V)\bigr), \label{eq:math-pade-monoid-p} \\
    Q_{\mathcal{R}_1\mathbin{+\!+}\mathcal{R}_2}(U, V) &= Q_{\mathcal{R}_2}\bigl(P_{\mathcal{R}_1}(U, V),\; Q_{\mathcal{R}_1}(U, V)\bigr), \label{eq:math-pade-monoid-q}
  \end{align}
  а сумма и разность числителя и знаменателя на любом этаже $\mathcal{R}$ факторизуются через мультипликативные инварианты $\Pi_+(\mathcal{R}) = \prod_{r\in\mathcal{R}} r$ и $\Pi_-(\mathcal{R}) = \prod_{r\in\mathcal{R}} (2 - r)$:
  \begin{align}
    P_{\mathcal{R}}(U, V) + Q_{\mathcal{R}}(U, V) &\equiv \Pi_+(\mathcal{R})\,(U + V), \label{eq:math-pade-sum-factor} \\
    P_{\mathcal{R}}(U, V) - Q_{\mathcal{R}}(U, V) &\equiv \Pi_-(\mathcal{R})\,(U - V) \;=\; \Pi_-(\mathcal{R})\,C. \label{eq:math-pade-diff-factor}
  \end{align}
\item \textbf{Бескорневое кубическое сжатие Паде, точное сокращение $X^3$ на всех этажах $\mathcal{R}$ и закон ухода на соседний этаж:}
  для двух пар кубов $(U, V) = (\widetilde{A}X^3, \widetilde{B}Y^3)$ ($C = U - V$) и $(U', V') = (\widetilde{A}X'^3, \widetilde{B}Y'^3)$ ($C' = U' - V'$), кросс-валов $a = Y'X$, $b = X'Y$, линейной формы этажа $I_{\mathcal{R}} = a P_{\mathcal{R}} - b Q_{\mathcal{R}}$ и строго положительной квадратичной формы $\Phi_3(u, v) = u^2 + u v + v^2$ на каждом этаже $\mathcal{R}$ выполняются тождество бескорневого кубического сжатия и тождество сокращения $X^3$:
  \begin{align}
    U\,I_{\mathcal{R}}\,\Phi_3(a P_{\mathcal{R}}, b Q_{\mathcal{R}}) &\equiv a^3\bigl(U P_{\mathcal{R}}^3 - V Q_{\mathcal{R}}^3\bigr) - (U b^3 - V a^3)\,Q_{\mathcal{R}}^3, \label{eq:math-arbitrary-floor-cubic-pade} \\
    \widetilde{A}\,I_{\mathcal{R}}\,\Phi_3(a P_{\mathcal{R}}, b Q_{\mathcal{R}}) &\equiv Y'^3\bigl(U P_{\mathcal{R}}^3 - V Q_{\mathcal{R}}^3\bigr) - C'\,Y^3\,Q_{\mathcal{R}}^3 \qquad (X^3 \ge 1), \label{eq:math-arbitrary-floor-x3-cancel} \\
    U P_{[3]}^3 - V Q_{[3]}^3 &\equiv C^3\,(U + V), \label{eq:math-base-floor3-collapse} \\
    \det\mathcal{W}_{2,3}(U, V) &:= \widehat{P}_1\widehat{Q}_2 - \widehat{Q}_1\widehat{P}_2 \;\equiv\; (U - V)^3\bigl(U^2 + 3UV + V^2\bigr). \label{eq:math-cross-floor-wronskian23}
  \end{align}
  При переходе с этажа $\mathcal{R}$ на соседний этаж $\mathcal{R}' = \mathcal{R} \mathbin{+\!+} [k_{\mathrm{shift}} + 1]$ линейная форма удовлетворяет разностному закону
  \begin{equation}
    \label{eq:math-adjacent-floor-escape}
    I_{\mathcal{R}'} - (k_{\mathrm{shift}} + 1)\,I_{\mathcal{R}} \;\equiv\; -k_{\mathrm{shift}}\,(a + b)\,\Pi_-(\mathcal{R})\,C,
  \end{equation}
  откуда при $k_{\mathrm{shift}}(a+b)\Pi_-(\mathcal{R})C \ne 0$ формы $I_{\mathcal{R}}$ и $I_{\mathcal{R}'}$ не могут обратиться в нуль одновременно ($\neg(I_{\mathcal{R}} = 0 \land I_{\mathcal{R}'} = 0)$ и $|I_{\mathcal{R}}| + |I_{\mathcal{R}'}| \ge 1$).
\end{enumerate}
\end{theorem}

\begin{proof}
\leanok
Уравнение конуса~\eqref{eq:math-sl2-isotropic-cone} следует из $\det\widetilde{P} = (S+X)(S-X) - (Y+T)(Y-T) = S^2 - X^2 - Y^2 + T^2 = 1$ при $S = \frac{1}{2}\operatorname{Tr}\widetilde{P}$. Ассоциативность моноида~\eqref{eq:math-pade-monoid-p}--\eqref{eq:math-pade-monoid-q} и факторизации~\eqref{eq:math-pade-sum-factor}--\eqref{eq:math-pade-diff-factor} доказываются структурной индукцией по списку $\mathcal{R}$. В пункте~(3) разность $(a P_{\mathcal{R}})^3 - (b Q_{\mathcal{R}})^3 = I_{\mathcal{R}}\,\Phi_3(a P_{\mathcal{R}}, b Q_{\mathcal{R}})$ после умножения на $U$ и выделения $a^3(U P_{\mathcal{R}}^3 - V Q_{\mathcal{R}}^3)$ даёт~\eqref{eq:math-arbitrary-floor-cubic-pade}; подстановка $U = \widetilde{A}X^3$, $V = \widetilde{B}Y^3$, $a^3 = Y'^3 X^3$, $b^3 = X'^3 Y^3$ выносит общий множитель $X^3$, сокращение которого даёт~\eqref{eq:math-arbitrary-floor-x3-cancel}. Наконец, подстановка $P_{\mathcal{R}'} = P_{\mathcal{R}} + k_{\mathrm{shift}}Q_{\mathcal{R}}$ и $Q_{\mathcal{R}'} = k_{\mathrm{shift}}P_{\mathcal{R}} + Q_{\mathcal{R}}$ даёт $I_{\mathcal{R}'} - (k_{\mathrm{shift}}+1)I_{\mathcal{R}} = -k_{\mathrm{shift}}(a+b)(P_{\mathcal{R}} - Q_{\mathcal{R}}) = -k_{\mathrm{shift}}(a+b)\Pi_-(\mathcal{R})C$, что доказывает~\eqref{eq:math-adjacent-floor-escape}.
\end{proof}

% ============================================================================
\chapter{Трёхшахтные планетарные дифференциалы \texorpdfstring{$\mathcal{A} + \mathcal{B} = \mathcal{C}$}{A + B = C}, вронскиан Мейсона--Стотерса и проекторы Лагранжа--Вандермонда}
\label{ch:math-planetary-differentials}
% ============================================================================

\section{Ортогональный репер трёхшахтного дифференциала \texorpdfstring{$\mathcal{F}_3$}{F3}, матрица моментов Ганкеля и целочисленный вронскиан Мейсона--Стотерса}

Любой трёхшахтный аддитивный баланс $A + B = C$ при факторизации валов $A = a_1 a_2 = \widetilde{A}X^3$, $B = b_1 b_2 = \widetilde{B}Y^3$, $C = c_1 c_2 = \widetilde{C}Z^3$ (с кубическим дефектом $S_3 = A + B - C$ и дефектами нижних этажей $S_1 = \widetilde{A}X + \widetilde{B}Y - \widetilde{C}Z$, $S_2 = \widetilde{A}X^2 + \widetilde{B}Y^2 - \widetilde{C}Z^2$) задаёт ортогональность вектора радиусов $\mathbf{r} = (a_1, b_1, c_1)$ и вектора угловых скоростей $\mathbf{\omega} = (a_2, b_2, -c_2)$ ($\mathbf{r}\cdot\mathbf{\omega} = A + B - C$).

\begin{theorem}[Ортогональный репер дифференциала $\mathcal{F}_3$, определитель Ганкеля--Вандермонда, целочисленный вронскиан Мейсона--Стотерса, проекторы Лагранжа--Вандермонда и стереографическое вложение в ротатор $\mathbb{Z}_{4Q}$]
\label{thm:math-abc-planetary-wronskian-full}
\lean{thm8_unified_3x3_pillai_abc_skeletal_matrix_and_radical_floor_embedding, thm9_pillai_abc_differential_bezout_and_pade_laws, thm12_reverse_pade_bezout_cubic_annihilation_and_radical_bound, thm13_integer_mason_stothers_cubic_wronskian_and_fourth_power_det, thm14_master_pillai_and_abc_bilateral_convergence_chains, thm14b_abc_hankel_vandermonde_floor_defect_trichotomy_and_bilateral_bezout_bracket, thm14c_abc_three_shaft_minor_closure_and_joint_beta_gamma_brackets, thm16_weierstrass_stereographic_abc_embedding_into_z512_rotator, thm30_unified_pillai_abc_collatz_planetary_differential_bridge}
\uses{thm:math-wronskian-weierstrass-quadratures, thm:math-pade-tower-cubic-compression}
\leanok
Для произвольных целых параметров $(\widetilde{A}, \widetilde{B}, \widetilde{C}, X, Y, Z, u_0, v_0, t)$ с весами валов $w_x = \widetilde{A}X$, $w_y = \widetilde{B}Y$, $w_z = \widetilde{C}Z$, полным радикал-скелетом $R_{\mathrm{full}} = (\widetilde{A}\widetilde{B}\widetilde{C})(X Y Z)$ и определителем Безу $d_{\mathrm{Bez}} = w_x v_0 - w_y u_0$ выполняются следующие безусловные тождества:
\begin{enumerate}[label=\textbf{(\arabic*)},leftmargin=*]
\item \textbf{Тождество Лагранжа--Бине--Коши репера $\mathcal{F}_3$ и факторизация определителя Ганкеля $3\times 3$:}
  для миноров $M_{BC} = b_1 c_2 + c_1 b_2$, $M_{CA} = c_1 a_2 + a_1 c_2$, $S_{AB} = a_1 b_2 - b_1 a_2$, $M_{AB} = a_1 b_2 + b_1 a_2$ и матрицы моментов Ганкеля $\mathcal{H}_3 = \bigl(S_{i+j-1}\bigr)_{i,j=0}^2$ ($S_k = \widetilde{A}X^k + \widetilde{B}Y^k - \widetilde{C}Z^k$) выполняются:
  \begin{align}
    M_{BC}^2 + M_{CA}^2 + S_{AB}^2 + (a_1 a_2 + b_1 b_2 - c_1 c_2)^2 &\equiv (a_1^2 + b_1^2 + c_1^2)(a_2^2 + b_2^2 + c_2^2), \label{eq:math-frame-f3-lagrange} \\
    (A + B)^2 + S_{AB}^2 &\equiv (A - B)^2 + M_{AB}^2 \;=\; (a_1^2 + b_1^2)(a_2^2 + b_2^2), \label{eq:math-brahmagupta-shafts} \\
    -\det\mathcal{H}_3 &\equiv \widetilde{A}\widetilde{B}\widetilde{C}\,(X - Y)^2(Y - Z)^2(Z - X)^2. \label{eq:math-hankel-vandermonde-det}
  \end{align}
\item \textbf{Калибровочно-инвариантный кубический вронскиан Мейсона--Стотерса и определитель четвёртых степеней:}
  для калибровочно-сдвинутого семейства $(\alpha(t), \beta(t)) = \bigl(u_0 Z^2 + t w_x,\ v_0 Z^2 + t w_y\bigr)$ и квадратичных валов $A^*(t) = \alpha(t) X^2$, $B^*(t) = \beta(t) Y^2$ кубический вронскиан $W^*(t) := A B^*(t) - B A^*(t)$ не зависит от калибровочного параметра $t$ ($\Delta_t W^*(t) \equiv 0$, сохранение Нётер) и удовлетворяет тождествам:
  \begin{align}
    w_x\,\beta(t) - w_y\,\alpha(t) &\equiv d_{\mathrm{Bez}}\,Z^2, \label{eq:math-ms-minor-z2} \\
    W^*(t) := A\,B^*(t) - B\,A^*(t) &\equiv d_{\mathrm{Bez}}\,(X^2 Y^2 Z^2), \label{eq:math-ms-cubic-wronskian} \\
    d_{\mathrm{Bez}}\,(A B C) &\equiv R_{\mathrm{full}}\,W^*(t), \label{eq:math-volume-wronskian-equiv} \\
    -\det\mathcal{W}_{\mathrm{MS}} &\equiv d_{\mathrm{Bez}}\,\bigl(X^4 + Y^4 + Z^4\bigr) + \bigl(Y^2 u_0 - X^2 v_0\bigr)\,S_3. \label{eq:math-ms-fourth-power-det}
  \end{align}
\item \textbf{Проекторы Лагранжа--Вандермонда по дефектам этажей $(S_1, S_2, S_3)$ и трёхшахтная система коллапса объёма:}
  на любых этажах $(S_1, S_2, S_3)$ выполняются точные проекционные тождества Лагранжа--Вандермонда и трёхшахтные уравнения дефектов миноров $d_{\mathrm{Min},Z} = Z^2 - (w_x\beta - w_y\alpha)$, $d_{\mathrm{Min},Y} = Y^2 - (w_x\gamma - w_z\alpha)$, $d_{\mathrm{Min},X} = X^2 - (w_z\beta - w_y\gamma)$:
  \begin{align}
    C\,(Z - X)(Z - Y) &\equiv Z^2\,\bigl(S_2(X + Y) - S_1 X Y - S_3\bigr), \label{eq:math-vandermonde-proj-c} \\
    A\,(X - Y)(X - Z) &\equiv X^2\,\bigl(S_1 Y Z - S_2(Y + Z) + S_3\bigr), \label{eq:math-vandermonde-proj-a} \\
    B\,(Y - X)(Y - Z) &\equiv Y^2\,\bigl(S_1 Z X - S_2(Z + X) + S_3\bigr), \label{eq:math-vandermonde-proj-b} \\
    w_y\,C + (\widetilde{B}^2\widetilde{C})(Y^2 Z)\,\alpha - R_{\mathrm{full}}\,\beta &\equiv (\widetilde{B}\widetilde{C}Y Z)\,d_{\mathrm{Min},Z}, \label{eq:math-three-shaft-collapse-c} \\
    w_z\,B + (\widetilde{B}\widetilde{C}^2)(Y Z^2)\,\alpha - R_{\mathrm{full}}\,\gamma &\equiv (\widetilde{B}\widetilde{C}Y Z)\,d_{\mathrm{Min},Y}, \label{eq:math-three-shaft-collapse-b} \\
    w_y w_z\,A - R_{\mathrm{full}}(w_z\beta - w_y\gamma) &\equiv R_{\mathrm{full}}\,d_{\mathrm{Min},X}. \label{eq:math-three-shaft-collapse-a}
  \end{align}
\item \textbf{Стереографическое вложение трёхшахтного дифференциала $\mathcal{C} := \mathcal{A} + \mathcal{B}$ в ротатор $\mathbb{Z}_{4Q}$ и двусторонняя целочисленная вилка захвата ячейки:}
  для любого четверть-периода $Q \in \mathbb{Z}$, фазы $\varphi_0 \in \mathbb{Z}$ ($\mathcal{C}_{\mathrm{rot}} = Q - \varphi_0,\ \mathcal{S}_{\mathrm{rot}} = \varphi_0$), резидуала ротатора $I_{\mathrm{rot}} = \mathcal{B}\mathcal{C}_{\mathrm{rot}} - \mathcal{A}\mathcal{S}_{\mathrm{rot}}$ и определителя Безу $g = \mathcal{A}v - \mathcal{B}u$ ($w = u + v$, $\Delta_{\mathrm{rot}} = \mathcal{C}_{\mathrm{rot}}v - \mathcal{S}_{\mathrm{rot}}u$) выполняются:
  \begin{align}
    Q\,\mathcal{A} &\equiv \mathcal{C}\,\mathcal{C}_{\mathrm{rot}} - I_{\mathrm{rot}}, \qquad
    Q\,\mathcal{B} \;\equiv\; \mathcal{C}\,\mathcal{S}_{\mathrm{rot}} + I_{\mathrm{rot}}, \label{eq:math-rotator-ab-split} \\
    \mathcal{C}\,\Delta_{\mathrm{rot}} - w\,I_{\mathrm{rot}} &\equiv g\,Q, \label{eq:math-rotator-bezout-det}
  \end{align}
  причём в главной ячейке захвата $0 \le I_{\mathrm{rot}} < \mathcal{C}$ произведения $Q\mathcal{A}$ и $Q\mathcal{B}$ зажаты в двусторонней целочисленной вилке $\mathcal{C}\mathcal{C}_{\mathrm{rot}} - \mathcal{C} < Q\mathcal{A} \le \mathcal{C}\mathcal{C}_{\mathrm{rot}}$ и $\mathcal{C}\mathcal{S}_{\mathrm{rot}} \le Q\mathcal{B} < \mathcal{C}\mathcal{S}_{\mathrm{rot}} + \mathcal{C}$.
\end{enumerate}
\end{theorem}

\begin{proof}
\leanok
Все полиномиальные тождества~\eqref{eq:math-frame-f3-lagrange}--\eqref{eq:math-rotator-bezout-det} следуют из раскрытия скобок в коммутативном кольце $\mathbb{Z}$.
\end{proof}

% ============================================================================
\chapter{Алгебра многовинтовых планетарных каскадов и встречно-волновая суперпозиция на сфере Римана \texorpdfstring{$\mathbb{P}^1$}{P1}}
\label{ch:math-multiscrew-and-holonomy}
% ============================================================================

\section{Моноид многовинтовых планетарных каскадов, экстремальная перестановочная вилка и вариационная монотонность}

Рассмотрим произвольный каскад линейных передач (винтов) $L = [(P_1, Q_1, D_1), \dots, (P_M, Q_M, D_M)] \in \mathrm{List}\,(\mathbb{Z}\times\mathbb{Z}\times\mathbb{Z})$, где каждая ступень связывает входной и выходной валы соотношением $P_k n_k - Q_k n_{k-1} = D_k$. Индуктивные операторы полного передаточного числа $P_{\mathrm{tot}}(L) = \prod P_k$, $Q_{\mathrm{tot}}(L) = \prod Q_k$, модуля зазора $C_{\mathrm{tot}}(L) = P_{\mathrm{tot}}(L) - Q_{\mathrm{tot}}(L)$ и накопленного дрейфа $D_{\mathrm{tot}}(L)$:
\begin{equation}
  \label{eq:math-multiscrew-d-rec}
  D_{\mathrm{tot}}\bigl((p, q, d) :: L_{\mathrm{rest}}\bigr)
  \;:=\;
  Q_{\mathrm{tot}}(L_{\mathrm{rest}})\cdot d + p\cdot D_{\mathrm{tot}}(L_{\mathrm{rest}})
\end{equation}
образуют полупрямой матричный моноид.

\begin{theorem}[Моноид многовинтовых каскадов, встречное сшивание Безу, экстремальная вилка бинарных слов шестерён $3^B - 2^B \le R_w \le 2^{A-B}(3^B - 2^B)$ и вариационная монотонность транспозиций]
\label{thm:math-multiscrew-cascade-extremal-envelope}
\lean{thm18_arbitrary_modulus_rotator_and_two_arc_trajectory_wronskian, thm24_collatz_counter_rotating_forward_backward_splice_and_bezout_lock, thm29_collatz_trailing_ones_burn_macro_step, thm31_collatz_binary_gear_word_extremal_drift_envelope_and_cycle_squeeze, thm32_collatz_countable_core_stage_monoid_dual_drift_and_cycle_gap, thm33_collatz_boundary_fixed_point_reflection_cycle_ratio_and_qmin_collapse, thm36_unified_extremum_principles_across_pillai_abc_and_collatz, thm37_collatz_resonant_gap_elimination_orbital_shift_and_qmin_qmax_sandwich}
\uses{thm:math-pade-tower-cubic-compression}
\leanok
Для любых списков винтов $L_1, L_2 \in \mathrm{List}\,(\mathbb{Z}\times\mathbb{Z}\times\mathbb{Z})$ и любого бинарного слова шестерён $w \in \mathrm{List}\,\mathrm{Bool}$ из нечётной ступени $\mathcal{G}_{\mathrm{odd}} = (2, 3, 1)$ и чётной ступени $\mathcal{G}_{\mathrm{even}} = (2, 1, 0)$ ($P_w = 2^A,\ Q_w = 3^B,\ P_{\mathrm{odd}} = 2^B,\ P_{\mathrm{even}} = 2^{A-B},\ R_w = D_{\mathrm{tot}}(L(w))$) выполняются следующие безусловные законы:
\begin{enumerate}[label=\textbf{(\arabic*)},leftmargin=*]
\item \textbf{Гомоморфизм моноида многовинтового каскада и закон двухдугового хордового вронскиана:}
  при конкатенации каскадов $L_1 \mathbin{+\!+} L_2$:
  \begin{align}
    P_{\mathrm{tot}}(L_1 \mathbin{+\!+} L_2) &= P_{\mathrm{tot}}(L_1)\,P_{\mathrm{tot}}(L_2), \qquad
    Q_{\mathrm{tot}}(L_1 \mathbin{+\!+} L_2) \;=\; Q_{\mathrm{tot}}(L_1)\,Q_{\mathrm{tot}}(L_2), \label{eq:math-multiscrew-pq-hom} \\
    D_{\mathrm{tot}}(L_1 \mathbin{+\!+} L_2) &= Q_{\mathrm{tot}}(L_2)\,D_{\mathrm{tot}}(L_1) + P_{\mathrm{tot}}(L_1)\,D_{\mathrm{tot}}(L_2), \label{eq:math-multiscrew-d-hom} \\
    C_{\mathrm{tot}}(L_1 \mathbin{+\!+} L_2) &= P_{\mathrm{tot}}(L_2)\,C_{\mathrm{tot}}(L_1) + Q_{\mathrm{tot}}(L_1)\,C_{\mathrm{tot}}(L_2). \label{eq:math-multiscrew-c-hom}
  \end{align}
  На любой замкнутой двухдуговой траектории ($n_{\mathrm{final}} = n_{\mathrm{start}}$, $D_1 = P_1 n_{\mathrm{mid}} - Q_1 n_{\mathrm{start}}$, $D_2 = P_2 n_{\mathrm{start}} - Q_2 n_{\mathrm{mid}}$) хордовый сдвиг удовлетворяет тождеству вронскиана $(P_1 P_2 - Q_1 Q_2)(n_{\mathrm{mid}} - n_{\mathrm{start}}) \equiv (P_2 - Q_2)D_1 - (P_1 - Q_1)D_2$.
\item \textbf{Экстремальная перестановочная вилка дрейфа по группе $S_A$ и строгая положительность вариационной производной транспозиции:}
  для любого бинарного слова шестерён $w \in \mathrm{List}\,\mathrm{Bool}$ выполняются $P_{\mathrm{odd}} P_{\mathrm{even}} = P_w$, $1 \le P_{\mathrm{odd}} \le Q_w$ и двусторонняя перестановочно-инвариантная оценка
  \begin{equation}
    \label{eq:math-binary-word-extremal-envelope}
    Q_w - P_{\mathrm{odd}} \;\le\; R_w \;\le\; P_{\mathrm{even}}\,(Q_w - P_{\mathrm{odd}}) \;=\; P_{\mathrm{even}} Q_w - P_w,
  \end{equation}
  а элементарная транспозиция соседних нечётной и чётной ступеней $[1, 0] \mapsto [0, 1]$ внутри $w_1 \mathbin{+\!+} [\dots] \mathbin{+\!+} w_2$ имеет строго положительную вариационную производную:
  \begin{equation}
    \label{eq:math-swap-variational-derivative}
    R_{w_1 \mathbin{+\!+} [0, 1] \mathbin{+\!+} w_2} - R_{w_1 \mathbin{+\!+} [1, 0] \mathbin{+\!+} w_2}
    \;\equiv\;
    P_{\mathrm{tot}}(L(w_1))\,Q_{\mathrm{tot}}(L(w_2)) \;\ge\; 1.
  \end{equation}
\item \textbf{Сохранение двойственного дрейфа нечётного ядра, мультипликативная вилка цикла и отражение граничных неподвижных точек:}
  для списка макроступеней нечётного ядра $L$ с верхним мультипликативным пакетом $U_L = \prod_{(P,Q,p)\in L} (Q + p - 1)$, прямым дрейфом $D_L = P_L q_f - Q_L q_0$ и двойственным верхним дефектом $E_L = U_L q_0 - P_L q_f$ выполняется закон сохранения двойственного дрейфа
  \begin{equation}
    \label{eq:math-dual-core-drift-conservation}
    D_L + E_L \;\equiv\; (U_L - Q_L)\,q_0,
  \end{equation}
  на любом замкнутом цикле ($q_f = q_0 \ge 1,\ D_L \ge 1,\ E_L \ge 0$) выполняются $(U_L - P_L)D_L = (P_L - Q_L)E_L$ и $Q_L < P_L \le U_L$, а граничные неподвижные точки чисто нечётного нижнего края $n_\infty^{(-)} = -1$ и чисто чётного верхнего края $n_\infty^{(+)} = 0$ при отражении дают стационарный корень $n^* = 2 n_\infty^{(+)} - n_\infty^{(-)} = +1$.
\end{enumerate}
\end{theorem}

\begin{proof}
\leanok
Гомоморфизмы~\eqref{eq:math-multiscrew-pq-hom}--\eqref{eq:math-multiscrew-c-hom} доказываются структурной индукцией по списку $L_1$. Оценка~\eqref{eq:math-binary-word-extremal-envelope} доказывается индукцией по списку $w \in \mathrm{List}\,\mathrm{Bool}$ с разбором чётной и нечётной ступени, а формула вариационной производной транспозиции~\eqref{eq:math-swap-variational-derivative} следует из двукратного применения~\eqref{eq:math-multiscrew-d-hom}.
\end{proof}

\section{Векторная суперпозиция встречных волн, тройная дихотомия \texorpdfstring{$M$}{M}-этажного каскада и модельно-теоретическая независимость}

\begin{theorem}[Тождественное сокращение чётной матрицы Грама $\mathrm{MainAmount}_{\mathrm{vec}}$ при векторной суперпозиции встречных волн, тройная дихотомия каскада и свидетели модельно-теоретической независимости]
\label{thm:math-counterwave-dichotomy-and-independence}
\lean{ocean_true_vector_superposition_norm_imbalance_and_mainamount_cancellation, counterwave_cascade_orthogonality_dichotomy_and_independence_models}
\uses{thm:math-pade-tower-cubic-compression, thm:math-multiscrew-cascade-extremal-envelope}
\leanok
Пусть $\mathcal{O}_{\mathrm{vec}} = \{(c_k, s_k, \alpha_k, \beta_k)\}_{k=1}^N \in \mathrm{List}\,(\mathbb{Z}^4)$ --- произвольный конечный ансамбль двумерных волновых мод с суммарным базовым фазором $(C_\Sigma, S_\Sigma) = \bigl(\sum_{k=1}^N c_k,\ \sum_{k=1}^N s_k\bigr)$, суммарным вектором радиальной производной $(A_\Sigma, B_\Sigma) = \bigl(\sum_{k=1}^N \alpha_k,\ \sum_{k=1}^N \beta_k\bigr)$, вектором прямой волны $W_+(\delta) = (C_\Sigma + A_\Sigma\delta,\ S_\Sigma + B_\Sigma\delta)$ и вектором встречной волны $W_-(\delta) = (C_\Sigma - A_\Sigma\delta,\ S_\Sigma - B_\Sigma\delta)$ при сдвиге $\delta \in \mathbb{Z}$. Тогда:
\begin{enumerate}[label=\textbf{(\arabic*)},leftmargin=*]
\item \textbf{Тождественное сокращение полной матрицы Грама $\mathrm{MainAmount}_{\mathrm{vec}}(\delta)$ со всеми $N^2$ перекрёстными слагаемыми:}
  квадраты норм суммарных векторов $U_{\mathrm{vec}}(\delta) = |W_+(\delta)|^2$ и $V_{\mathrm{vec}}(\delta) = |W_-(\delta)|^2$ удовлетворяют:
  \begin{align}
    U_{\mathrm{vec}}(\delta) + V_{\mathrm{vec}}(\delta) &\equiv 2\,\mathrm{MainAmount}_{\mathrm{vec}}(\delta)
      \;=\; 2\,\Bigl((C_\Sigma^2 + S_\Sigma^2) + \delta^2(A_\Sigma^2 + B_\Sigma^2)\Bigr), \label{eq:math-vec-ocean-sum} \\
    U_{\mathrm{vec}}(\delta) - V_{\mathrm{vec}}(\delta) &\equiv 4\,\bigl(C_\Sigma A_\Sigma + S_\Sigma B_\Sigma\bigr)\,\delta. \label{eq:math-vec-ocean-diff}
  \end{align}
\item \textbf{Тройная дихотомия $M$-этажного каскада $\mathrm{SL}_2(\mathbb{Z})$:}
  после прохождения пары $(U_{\mathrm{vec}}(\delta), V_{\mathrm{vec}}(\delta))$ через произвольный $M$-этажный каскад $R \in \mathrm{List}\,\mathbb{Z}$ выходные амплитуды $P_R = \operatorname{minkowskiCascadeP}(R, U_{\mathrm{vec}}(\delta), V_{\mathrm{vec}}(\delta))$ и $Q_R = \operatorname{minkowskiCascadeQ}(R, U_{\mathrm{vec}}(\delta), V_{\mathrm{vec}}(\delta))$ удовлетворяют
  \begin{equation}
    \label{eq:math-cascade-diff-trichotomy}
    P_R - Q_R \;\equiv\; 4\,\Pi_-(R)\,\bigl(C_\Sigma A_\Sigma + S_\Sigma B_\Sigma\bigr)\,\delta,
  \end{equation}
  а при $q > 0$, $P_R > 0$, $Q_R > 0$ обращение в нуль интенсивности деструктивной интерференции $I_q(P_R, Q_R, \mathcal{C}_q(\Delta\varphi)) = q(P_R - Q_R)^2 + 2 P_R Q_R(q + \mathcal{C}_q(\Delta\varphi))$ эквивалентно тройной дихотомии:
  \begin{equation}
    \label{eq:math-counterwave-zero-iff-trichotomy}
    I_q\bigl(P_R,\, Q_R,\, \mathcal{C}_q(\Delta\varphi)\bigr) = 0
    \;\Longleftrightarrow\;
    \Bigl(\Pi_-(R) = 0 \;\lor\; C_\Sigma A_\Sigma + S_\Sigma B_\Sigma = 0 \;\lor\; \delta = 0\Bigr)
    \;\land\; \operatorname{pauliPhaseDistNat}(q, \Delta\varphi) = 0.
  \end{equation}
\item \textbf{Две конструктивные модели и теорема о модельно-теоретической независимости $\delta = 0$ от одних лишь встречно-волновых аксиом:}
  \begin{itemize}[noitemsep]
  \item \textbf{Модель~1 (коллинеарная модель $\mathcal{O}_1 = [(1, 0, 1, 0)]$):} $(C_\Sigma, S_\Sigma) = (1, 0)$, $(A_\Sigma, B_\Sigma) = (1, 0)$, скалярное произведение $C_\Sigma A_\Sigma + S_\Sigma B_\Sigma = 1 \ne 0$, откуда при $\Pi_-(R) \ne 0$ равенство $P_R = Q_R$ эквивалентно $\delta = 0$;
  \item \textbf{Модель~2 (ортогональная модель $\mathcal{O}_2 = [(1, 1, -1, 1)]$):} $(C_\Sigma, S_\Sigma) = (1, 1) \perp (A_\Sigma, B_\Sigma) = (-1, 1)$, скалярное произведение $C_\Sigma A_\Sigma + S_\Sigma B_\Sigma = 1\cdot(-1) + 1\cdot 1 = 0$, и на ненулевом сдвиге $\delta = 1 \ne 0$ выполняются $U_{\mathrm{vec}}(1) = 4 > 0$, $V_{\mathrm{vec}}(1) = 4 > 0$, $P_R = Q_R$ и $I_q(P_R, Q_R, \mathcal{C}_q(2q)) = 0$.
  \end{itemize}
  Следовательно, утверждение $\delta = 0$ модельно-теоретически независимо от аксиом встречно-волновой симметрии и каскада $\mathrm{SL}_2(\mathbb{Z})$ без дополнительного аналитического закона, исключающего ортогональность $(C_\Sigma, S_\Sigma) \perp (A_\Sigma, B_\Sigma)$.
\end{enumerate}
\end{theorem}

\begin{proof}
\leanok
Тождества~\eqref{eq:math-vec-ocean-sum}--\eqref{eq:math-counterwave-zero-iff-trichotomy} и свойства моделей $\mathcal{O}_1, \mathcal{O}_2$ следуют из билинейности скалярного произведения и факторизации разности башни Паде.
\end{proof}

\section{Конструктивное построение \texorpdfstring{$\mathbb{R}$ и $\mathbb{C}$}{R and C}, биективный изоморфизм \texorpdfstring{$L^1(\mathbb{Z}_{4Q}) \leftrightarrow S^1(\mathbb{C})$}{L1 <-> S1(C)} и бритва Оккама--Мальцева}

\begin{theorem}[Конструктивное построение полей $\mathbb{R}$ и $\mathbb{C}$ через вилки Штерна--Броко--$\mathrm{SL}_2(\mathbb{Z})$, биективный стереографический изоморфизм с $S^1(\mathbb{C})$ и теорема Мальцева об изоморфизме элементарной теории]
\label{thm:math-constructive-reals-s1-malcev-limit}
\lean{thm17_model_theoretic_termination_kernel_annihilation_and_godel_tarski_oracle_necessity, thm41_constructive_projective_real_intervals_stern_brocot_and_s1_isomorphism}
\uses{thm:math-wronskian-weierstrass-quadratures, thm:math-counterwave-dichotomy-and-independence}
\leanok
Поля $\mathbb{Q}$, $\mathbb{R}$ и $\mathbb{C}$ (включая единичную окружность $S^1(\mathbb{C})$) строятся полностью конструктивно (\texttt{0 Classical.choice}) как проективный предел вложенных рациональных вилок и счётной башни носителей $(M_n, \Pi_{n\to 0})$. В силу изоморфизма фактор-системы по ядру гомоморфизма Мальцева $\ker\pi$ и совпадения элементарной теории первого порядка $\mathrm{Th}(M_{\mathrm{ext}}) \equiv \mathrm{Th}(M_{\mathrm{base}})$, по бритве Оккама не требуется вводить никаких отдельных искусственных семейств вроде <<$\mathbb{R}_{\mathrm{proj}}$>> или <<$\mathbb{C}_{\mathrm{proj}}$>> --- построенные объекты изоморфны обычным $\mathbb{Q}$, $\mathbb{R}$ и $\mathbb{C}$, которыми мы и оперируем напрямую:
\begin{enumerate}[label=\textbf{(\arabic*)},leftmargin=*]
\item \textbf{Конструктивное представление чисел из $\mathbb{R}$ через вилки Штерна--Броко и действие $\mathrm{SL}_2(\mathbb{Z})$:}
  каждая рациональная вилка $I = [p_L/q_L,\, p_R/q_R]$ ($q_L, q_R > 0$) характеризуется целочисленным вронскианом (определителем) $\Delta(p_L, q_L, p_R, q_R) := p_R q_L - p_L q_R \ge 0$. Деление вилки медиантой Штерна--Броко $(p_M, q_M) = (p_L + p_R,\ q_L + q_R)$ тождественно сохраняет вронскиан обеих дочерних вилок:
  \begin{equation}
    \label{eq:math-stern-brocot-wronskian}
    \Delta(p_L, q_L, p_M, q_M) \;=\; \Delta(p_M, q_M, p_R, q_R) \;\equiv\; \Delta(p_L, q_L, p_R, q_R),
  \end{equation}
  а под действием любой целочисленной матрицы уточнения $A = \begin{pmatrix} a & b \\ c & d \end{pmatrix} \in \mathrm{Mat}_{2\times 2}(\mathbb{Z})$ на концы вилки вронскиан умножается на определитель матрицы:
  \begin{equation}
    \label{eq:math-sl2-interval-wronskian}
    \Delta(p_L', q_L', p_R', q_R') \;\equiv\; (a d - b c)\,\Delta(p_L, q_L, p_R, q_R),
  \end{equation}
  при этом вдоль любого слова уточнений $w \in \mathrm{List}\,\mathbb{Z}$ свёрнутая вилка $[L_w, U_w] = \operatorname{bracketFold}(l_0, u_0, w)$ монотонно зажимает вещественное значение $X \in \mathbb{R}$ ($l_0 \le L_w \le X \le U_w \le u_0$, $0 \le U_w - L_w \le u_0 - l_0$) и однозначно определяет его при смыкании ($U_w - L_w = 0 \implies L_w = U_w = X$).
\item \textbf{Биективный стереографический изоморфизм $L^1(\mathbb{Z}_{4Q}) \leftrightarrow S^1(\mathbb{C})$ на всех масштабах $Q \in \mathbb{Z}$:}
  для любой точки ротатора $(C, S)$ с $C + S = Q$ прямой проективный подъём Вейерштрасса $(X, Y, R) = (C^2 - S^2,\ 2CS,\ C^2 + S^2)$ лежит на единичной окружности $X^2 + Y^2 = R^2$ в $\mathbb{C}$ ($R + Y = Q^2$) и биективно обращается без извлечения квадратных корней по линейным формулам:
  \begin{equation}
    \label{eq:math-weierstrass-bijection-inverse}
    Q^2 + X \;\equiv\; 2\,Q\,C, \qquad Q^2 - X \;\equiv\; 2\,Q\,S.
  \end{equation}
\item \textbf{Теорема Мальцева об изоморфизме элементарной теории $\mathrm{Th}(M_{\mathrm{ext}}) \equiv \mathrm{Th}(M_{\mathrm{base}})$ (бритва Оккама для $\mathbb{R}$ и $\mathbb{C}$):}
  для любой счётной проективной башни носителей $M_{k+1} \xrightarrow{\pi_k} M_k$ проекция $\Pi_{n \to 0} : M_n \to M_0$ сохраняет все согласованные наблюдаемые ($\operatorname{obs}_0(\Pi_{n \to 0}(x_n)) = \operatorname{obs}_n(x_n)$); для любого сюръективного гомоморфизма алгебраических систем Мальцева $\pi : M_{\mathrm{ext}} \twoheadrightarrow M_{\mathrm{base}}$ с сечением $\mathrm{sect} : M_{\mathrm{base}} \hookrightarrow M_{\mathrm{ext}}$ ($\pi(\mathrm{sect}(m)) = m$) фактор по конгруэнции ядра $\ker\pi$ изоморфен $M_{\mathrm{base}}$, а элементарные теории первого порядка тождественно совпадают:
  \begin{equation}
    \label{eq:math-malcev-elementary-equivalence}
    \forall\,\phi \in \mathrm{Formula}, \qquad
    \Bigl(\forall\,y \in M_{\mathrm{ext}},\ M_{\mathrm{ext}}, y \models \phi\Bigr)
    \;\Longleftrightarrow\;
    \Bigl(\forall\,m \in M_{\mathrm{base}},\ M_{\mathrm{base}}, m \models \phi\Bigr).
  \end{equation}
\end{enumerate}
\end{theorem}

\begin{proof}
\leanok
Тождества~\eqref{eq:math-stern-brocot-wronskian}--\eqref{eq:math-weierstrass-bijection-inverse} проверяются раскрытием скобок в кольце $\mathbb{Z}$; монотонное сжатие вилки и стабилизация башни $\Pi_{n\to 0}$ доказываются структурной индукцией по слову $w$ и номеру этажа $n$, а эквивалентность~\eqref{eq:math-malcev-elementary-equivalence} следует из $\pi(\mathrm{sect}(m)) = m$.
\end{proof}
